\begin{abstract}
\noindent
Quirks are heavy particles charged under a new confining gauge group with no light matter, so that a produced pair stays
connected by a macroscopic string of tension $\sigma=\Lam^2$. For colour-neutral, electroweak-charged quirks with
$\Lam\sim0.1$--$100\keV$ the pair oscillates with an amplitude between metres and micrometres and leaves tracker hits that are
non-helical and nearly coplanar with the primary vertex. We study two questions for an HL-LHC, CMS Phase-2-like tracker
($B_0=3.8\tesla$, $\sqrt s=14\TeV$, pile-up 200, $3\abinv$), using a vector-like lepton-like benchmark $E=(1,1)_{-1}$
with $N_{\rm IC}=2$ and its scalar partner.

\emph{(P1, truth level.)} We derive how the solenoid field thickens the quirk plane. In the pair rest frame the field
acts as a weak electric field that makes the two ends miss each other by $\bodd$ at every second crossing; the string's
transverse pull $\sigma\bodd/|x|$ then rotates the oscillation axis by an amount carrying $\ln(\lmax/\bodd)$. The
accumulated plane thickness follows $\DsB\propto\epsB[\ln(1/\epsB)+C]$ with $\epsB=qB/\Lam^2$, i.e. a local
logarithmic slope $c=1/[\ln(1/\epsB)+C]\simeq0.104$ at $\Lam=1\keV$, $\mQ=500\GeV$, derived analytically and
confirmed by blind numerics at $1\keV$ ($0.105\pm0.002$). In a $2.2\times10^5$-event truth-level simulation with the
full production kinematics the law holds per event (paired agreement $\le1\%$ at $1\keV$, thickness within $3.4\%$ for
the fermion and $5.3\%$ for the scalar), but the population median of $c$ lies $+2$--$3\%$ above the prediction at
$1\keV$, above a pre-set $1.5\%$ validation gate, which therefore failed; we read this as an estimator effect of the
per-event scatter (estimate). The law holds for median $\lmax/\RIT\lesssim0.6$ and $\kB\lesssim4$; at $3$--$5\keV$ a
larger, $\kB$-dependent excess is measured whose attribution to the centre-of-mass magnetic field is an estimate; one
prediction (a straight-string thickness ratio at $\gst=2$) fails at $\approx3\sigma$. Multiple scattering
sets a measured floor of $10$--$28\um$ and dominates above $\Lam\approx1.2$--$1.9\keV$.

\emph{(P2, detector level.)} Truth-level quirk hits are passed through an ACTS/Fatras model of the public Phase-2 tracker layout
with a pile-up-200 overlay normalised to the Tracker TDR occupancies, and searched for with a dedicated barrel plane
finder ($|\eta|<1$) under a working-point rule frozen before any signal efficiency was computed.
\textbf{No working point exists}: the lowest expected background is $1.2\times10^6$ events in $3\abinv$, 7 decades above the
budget, and beyond the width cuts it consists of zero-count bounds from finite simulated samples, so the result is limited
by simulation statistics and by the pile-up extrapolation. The pre-registered fallback gives expected limits at best 790
times the predicted cross section ($\mQ=200\GeV$, $\Lam=1.5\keV$; corrected for the $p_T(Q\bar Q)>200\GeV$ tail fraction,
omitted in an earlier version); nothing is excluded, and the existing bounds (FASER, D0,
recasts) are not extended. At pile-up 200 neither $\mQ$ nor $\Lam$ is measured: the time-of-flight estimator is a pair-mass
proxy with a pile-up-induced low cluster, and the $\Lam^{-2}$ pair-separation handle is lost. A data-driven pile-up
background estimate would be the decisive next step.
\end{abstract}
